\documentclass[12pt]{lettre} 
\usepackage[utf8]{inputenc} 
\usepackage[T1]{fontenc} 
\usepackage{lmodern} 
\usepackage{eurosym} 
%\usepackage[frenchb]{babel} 
\usepackage{numprint} 


\begin{document}


\begin{letter}{Monsieur le Préfet de la Meuse\\Préfecture de la Meuse\\55000 BAR LE DUC}




\address{Clément \textsc{Fevez}\\ Maire de Evres\\Place de la Halle\\55250 EVRES}
\lieu{Evres}
\telephone{03 29 92 05 97}
\email{evres@orange.fr}
\nofax



\def\concname{Objet :~} % On définit ici la commande 'objet' 
\conc{Demande de mise en place d'un système de vidéoprotection à Evres} 

\opening{Monsieur le Préfet,} 
J'ai l'honneur de vous soumettre un dossier de demande de mise en place d'un système de vidéoprotection au sein du village d'Evres.
Cette mise en place de 6 caméras repose sur la poursuite des finalités suivantes:


\begin{itemize}

\item Protection des bâtiments et installations publics et de leurs abords

\item Prévention des atteintes à la sécurité des personnes et des biens dans des lieux particulièrement exposés à des risques d'agression et de vol ou de trafic de stupéfiant

\item Prévention et constatation des infractions relatives à l'abandon d'ordures, de déchets, de matériaux ou d'autres objets

\end{itemize}
Au delà de ces finalités légales, le but poursuivi est de diminuer le sentiment d'impunité ressenti par les rôdeurs dans la commune.

En effet, il est fréquent à Evres que des habitants constatent la présence de véhicules au comportement suspect, ayant le plus souvent permis à des rôdeurs de venir reconnaître ou voler dans le village.

Ainsi, des vols de carburant ont plusieurs fois été signalés sur des tracteurs garés dans le village. Les habitants concernés, étant donné le préjudice et la certitude de ne jamais pouvoir le voir compenser, ne portent jamais plainte.
Deux fois récemment, des rôdeurs à pied ont été décelés et mis en fuite par des habitants. Lors d'un de ces événements, les rôdeurs avaient pris soin de détériorer le système d'éclairage communal dans les armoires électriques.
Très récemment, un rôdeur s'est introduit en pleine nuit au domicile d'un habitant en situation de handicap, et a été mis en fuite par l'habitant lui même.

Parfois, et malgré nos incitations à la prudence, ces événements provoquent des réactions manquant de tempérance de la part de nos habitants, ce que je peux comprendre. Par exemple, un habitant a récemment poursuivi une voiture suspecte dans laquelle était remonté un rôdeur jusqu'à un village voisin, se mettant en danger sur la route. Une autre fois, un habitant chasseur est sorti en pleine nuit et a tiré deux coups de feu en l'air pour faire fuir deux rôdeurs, qui avait été décelés par un voisin, ce dernier ayant donné l'alerte par un fil whatsapp.

En tant que maire, cette situation m'inquiète, et nous pensons avec mon conseil municipal que la mise en place d'un système de vidéoprotection, si elle ne supprimera pas tous les risques, permettra de diminuer la facilité et le sentiment d'impunité avec lesquels des voleurs viennent dans notre village pour y perpétrer des larcins. Par ailleurs, après la commission éventuelle d'une infraction, les images pourraient être utilisées par les forces de l'ordre pour faciliter les instructions.

Espérant que cette demande soit acceptée, je vous prie d'agréer, monsieur le Préfet, l'expression de mes meilleurs sentiments.


\end{letter} 

\end{document}
